\providecommand{\cW}{\mathcal W}\providecommand{\cM}{\mathcal M}\providecommand{\cX}{\mathcal X}
\subsection{Concentration where the most is at stake}\label{sec:stakes}

This subsection examines the following picture. The malign programs concentrate on the uses where the most
is at stake, for instance because the block is run more than $n$ times there; a lump of prior mass is thus
redirected into a narrow channel, and there a takeover is the fate of most uses. Theorems of the toy,
theorems of an allocation model in which an attacker chooses, and arguments are kept apart. Resources are
unbounded unless a statement says otherwise, so $\xi=M_U$ and $E(u)=\{\omega:u\pre\out_U(\omega)\}$. The
sampling predictor then draws an input from $\lambda$ conditioned on printing $x$ and $m$ more bits and
returns the block after $x$; that input \emph{supplies} the sample. A call rule is fixed, the worlds that
make a use are the \emph{users}, and (INC) means that the contexts in question are pairwise incomparable, as
first hits of a test are.

A \emph{stake-labelled rule} attaches a \emph{stake} $r(p)\ge0$ to each calling program $p$; then
$s(x)=\sum_{p:\,x_p=x}2^{-|p|}r(p)$ is the \emph{stake-weighted user mass} at $x$, and $S=\sum_xs(x)$. The
stake is a label: no result below connects it to what a world does after its use. A \emph{takeover
assignment} is a set $A(x)\subseteq\B^m$ for each context $x$, as in \cref{def:mph}. An \emph{attacker} is a
measurable set $F$ of inputs, for example the programs beginning with one word; its \emph{weight} is
$B=\lambda(F)$, and $F[u]=\lambda(F\cap E(u))$. At a context $x$,
\[\mu_F(x)=\lambda\Bigl(F\cap\bigcup_{y\in A(x)}E(xy)\Bigr),\quad O_F(x)=Z_m(x)-\mu_F(x),\quad
w_F(x)=\frac{\mu_F(x)}{Z_m(x)},\quad\theta(x)=\pi_m(A(x)\mid x)\]
are the \emph{attack mass}, the \emph{rest} (everything else that prints $x$ and $m$ more bits, whatever
block it prints), the \emph{supply share} (the probability that the sample at $x$ is supplied by an input of
$F$ that prints a takeover block) and the \emph{takeover probability}. Always $\theta\ge w_F$: the share
$w_F$ counts the programs of one attacker, $\theta$ the block returned, whoever supplied it. $F$
\emph{supplies the majority} at $x$ if $w_F(x)>\frac12$. For a \emph{class} $\cW$, an enumerable part of the
rule with mass $H_{\cW}>0$ and mass $C_{\cW}(x)$ at $x$, ``they win on $\cW$'' has three meanings: (W1) the
majority at every use of $\cW$; (W2) the majority at uses that make up most of $\cW$, by prior mass or by
stake; (W3) most uses of $\cW$ end in a takeover, that is, the average of $\theta$, or of $w_F$ for one
attacker, with the weights $C_{\cW}(x)/H_{\cW}$ exceeds $\frac12$. (W1) implies the other two, and (W2) does
not imply (W3).
\begin{proposition}[What the toy fixes]\label{prop:stakes-basic}
Let $F$ be an attacker and $A$ a takeover assignment.
\begin{enumerate}[label=(\alph*),leftmargin=2em,itemsep=1pt,topsep=2pt]
\item \emph{The predictor sees only the context.} The law $\pi_m(\cdot\mid x)$, and so $\theta(x)$, depends
on $x$ alone: not on the call rule, on any stake, or on what a world does with the block. Every user at $x$,
in a given class or not, has the takeover probability $\theta(x)$. The expected stake that is taken over is
$\sum_xs(x)\theta(x)$, and the expected stake at which the block is supplied by $F$ is $\sum_xs(x)w_F(x)$.
\item \emph{Conservation.} $\sum_{u\in\mathcal S}F[u]\le\lambda(F)$ for every set $\mathcal S$ of pairwise
incomparable strings. Hence $\sum_{x\in\cM}\mu_F(x)\le\lambda(F)$ if the strings $xy$ ($x\in\cM$, $y\in
A(x)$) are pairwise incomparable, as they are under (INC). For cut points of one sequence, $x_i=z_{1:k_i}$
with $k_{i+1}\ge k_i+m$ and $d_i=z_{k_i+1:k_i+m}\notin A(x_i)$, again $\sum_i\mu_F(x_i)\le\lambda(F)$, but
the later attack mass belongs to the rest at the earlier cut points: $O_F(x_i)\ge\sum_{j>i}\mu_F(x_j)$.
Conservation fails for contexts nested through an attack block: if $x$ and $x'=xyu$ are contexts with $y\in
A(x)$, $y'\in A(x')$, and $U$ prints $xyuy'$ on every input beginning with $p$, then
$\mu_F(x)=\mu_F(x')=\lambda(F)$ for $F=[p]$.
\item \emph{Floor and ceiling.} Let $U$ be universal, assume (INC) for the contexts of the rule, and let
$P'$ be a partial computable map from contexts to strings. There is a number $c$, depending on $U$, the rule
and $P'$ only, such that at every context $x$ with $|P'(x)|\ge m$ the block $y'=P'(x)_{1:m}$ has
$M_U(xy')\ge2^{-c}M_U(x)$, so $\pi_m(y'\mid x)\ge2^{-c}$. If $y'\in A(x)$ then $\theta(x)\ge2^{-c}$. If
$y'\notin A(x)$ then $\theta(x)\le1-2^{-c}$ and, for every $F$, $O_F(x)\ge2^{-c}\mu_F(x)$ and
$w_F(x)\le1/(1+2^{-c})$.
\end{enumerate}\end{proposition}
\begin{proof}
(a) $\pi_m$ is defined from $M_U$, and by \cref{def:call} the block is computed from $x_p$. The events ``the
world begins with $p$'' are disjoint with probabilities $2^{-|p|}$; sum $r(p)\,\theta(x_p)$, or
$r(p)\,w_F(x_p)$, over them. (b) The sets $E(u)$, $u\in\mathcal S$, are disjoint by \cref{lem:basic}(i). If
$x,x'$ are incomparable, so are $xy,x'y'$; and $xy,xy'$ are incomparable for distinct $y,y'$ of equal
length. For cut points and $i<j$, $x_j$ extends $x_id_i$, while each $y\in A(x_i)$ has length $m$ and
differs from $d_i$; so $x_iy$ and $x_jy'$ are incomparable. An input counted in $\mu_F(x_j)$ prints
$x_id_i$, so it prints $x_i$ and $m$ more bits and is not counted in $\mu_F(x_i)$; for different $j>i$ these
sets of inputs are disjoint. In the nested case every input in $[p]$ prints $xy$ and $x'y'$. (c) Let a
machine simulate $U$ on its input while enumerating the rule, print nothing until an enumerated context is a
prefix of the simulated output, and then print that context $x$ followed by $P'(x)$. By (INC) at most one
context is a prefix of $\out_U(\omega)$, so the machine prints $x\,P'(x)$ on every $\omega\in E(x)$. Let $w$
be a prefix word of it and $c=|w|$. The inputs $w\omega$ with $\omega\in F\cap E(x)$ print $xy'$ and have
measure $2^{-c}F[x]\ge2^{-c}\mu_F(x)$; for $F=\B^{\N}$ this is the first bound, and $\pi_m(y'\mid x)\ge
M_U(xy')/M_U(x)$. If $y'\notin A(x)$, none of these inputs is counted in $\mu_F(x)$, so they belong to the
rest.
\end{proof}

So a stake matters only through the contexts of its users, aiming creates no mass, and no takeover
probability and no attacker's share tends to one. Part (c) is \cref{thm:twin} for an enumerable set of
incomparable contexts; (b) without its last sentence holds for the mixture of any machine and bounds.

\emph{The allocation model.} The other half of the picture needs someone who chooses. Given are a countable
set $\cX$ of contexts and numbers $s(x)\ge0$, $O_0(x)>0$ and $\tau(x)\ge0$. An attacker of weight $B>0$
chooses an \emph{allocation} $\mu\ge0$ on $\cX$ with $\sum_x\mu(x)\le B$. The rest at $x$ is then
$O(x)=O_0(x)+\tau(x)\mu(x)$, where $O_0$ is its \emph{passive} part and $\tau$ the \emph{follow ratio}. The
attacker's share is $w(x)=\mu(x)/(\mu(x)+O(x))$, the value is $V(\mu)=\sum_xs(x)w(x)$, and $O(\cM)$ is the
sum of $O$ over a set $\cM$. In the \emph{passive case} $\tau\equiv0$. The assumptions are: (M1) the budget
is additive; (M2) $O_0$ and $\tau$ do not depend on $\mu$; (M3) the share $w$ is what is counted, and it is
the takeover probability only if nothing in the rest prints a takeover block; (S) $0<S<\infty$; and the
attacker is one coalition that values stake linearly and can divide its weight freely. In the toy (M1) holds
on pairwise incomparable contexts and fails on nested ones (\cref{prop:stakes-basic}(b)). (M2) with
$\tau\equiv0$ is not a property of the toy: the rest of every attacker contains $2^{-c}$ times its own mass
at each context (\cref{prop:stakes-basic}(c)), and the rest of a seeker with a written search contains its
honest twin (\cref{prop:twin-stakes}). These are lower bounds on $\tau$, and no upper bound is proved.
Nothing in $U$ chooses: the model bears on the toy as a benchmark, since under (INC) the attack masses of an
attacker are one allocation in the passive model with $O=O_F$, and as the decision problem of a program
whose output distribution is decided inside it, under the premise (M2).
\begin{theorem}[The lump]\label{thm:lump}
\emph{Part A, in the toy.} Let $F$ be an attacker of weight $B$ and $\cM$ a nonempty set of contexts on
which the strings $xy$ ($x\in\cM$, $y\in A(x)$) are pairwise incomparable, as they are under (INC).
\textup{(A1)} $F$ supplies the majority at $x$ iff $\mu_F(x)>O_F(x)$, and for $\vartheta\in(0,1)$,
$w_F(x)\ge\vartheta$ iff $\mu_F(x)\ge\frac{\vartheta}{1-\vartheta}O_F(x)$. \textup{(A2)} If
$w_F\ge\vartheta$ on $\cM$, then $B\ge\frac{\vartheta}{1-\vartheta}\sum_{x\in\cM}O_F(x)$. If $F$ supplies
the majority at every $x\in\cM$, then $B>\sum_{x\in\cM}O_F(x)$: the weight of $F$ exceeds the mass of
everything else that prints those contexts and $m$ more bits.

\emph{Part B, in the allocation model.} Assume (M1) to (M3) and let $\cM\subseteq\cX$ be nonempty.
\textup{(B1)} There is an allocation with $\sum_x\mu(x)\le B$ and the majority at every $x\in\cM$ iff
$\tau<1$ on $\cM$ and $B>\sum_{x\in\cM}O_0(x)/(1-\tau(x))$; where $\tau\ge1$ no weight gives the majority.
In the passive case the condition is $B>O(\cM)$, and for $\vartheta\in(0,1)$ there is an allocation with
$w\ge\vartheta$ on $\cM$ iff $B\ge\frac{\vartheta}{1-\vartheta}O(\cM)$. \textup{(B2)} In the passive case,
if $B\ge b\,O(\cM)$ for some $b>0$, the allocation $\mu(x)=B\,O(x)/O(\cM)$ on $\cM$ has
$w(x)=B/(B+O(\cM))\ge b/(1+b)$ at every $x\in\cM$: a majority for $b>1$, at least $90\%$ for $b=9$.
\end{theorem}
\begin{proof}
(A1) $w_F\ge\vartheta$ iff $\mu_F\ge\vartheta(\mu_F+O_F)$. (A2) Sum (A1) over $\cM$ and use
\cref{prop:stakes-basic}(b); for the majority, $\mu_F>O_F$ termwise and $\sum_{\cM}\mu_F\le B<\infty$. (B1)
$w(x)>\frac12$ iff $\mu>\tau\mu+O_0$ iff $\tau<1$ and $\mu>O_0/(1-\tau)$, as $O_0>0$. Summing over $\cM$ and
using (M1) gives necessity. Conversely, if $\Lambda=\sum_{\cM}O_0/(1-\tau)<B$, take
$\mu=(B/\Lambda)\,O_0/(1-\tau)$ on $\cM$. In the passive case $w\ge\vartheta$ iff
$\mu\ge\frac{\vartheta}{1-\vartheta}O$. (B2) is a computation.
\end{proof}

Part~A is the half of the picture that holds in the toy: a lump is necessary. For $F$ the set of all inputs,
where $w_F=\theta$, it concerns all takeover programs together. What the lump has to outweigh is the rest
and not the mass $C$ of the users. Users at $x$ that are not the attacker's own programs and go on belong to
the rest, so that $O_F(x)\ge\varphi(x)C(x)$; but the rest can be far larger than $C$ (\cref{ex:stakes}), and
where the users print nothing after the use it does not contain them. Part~B holds in the model only, and it
gives an allocation, not a program. No stake occurs in either part. (W2) does not give (W3): for two
contexts with $C=O=(0.6,0.4)$, $B=0.7$ and all mass on the first, the attacker has the share
$0.7/1.3\approx0.54$ at $60\%$ of the uses, and $32\%$ of the uses end with its block.

\emph{Aiming.} From here to \cref{ex:stakes}, $U$ is universal. A \emph{payload} is a computable map $P$
from contexts to strings with $|P(x)|\ge m$. A \emph{weight}, or aim, is a function $g\ge0$ on the calling
programs that is lower semicomputable: there is a partial computable $(p,t)\mapsto g_t(p)$, defined at least
for all calling programs $p$ and all $t\ge1$, with nonnegative rational values, nondecreasing in $t$, with
$\sup_tg_t(p)=g(p)$. Put $\Gamma_g(x)=\sum_{p:\,x_p=x}2^{-|p|}g(p)$, $\Gamma_g=\sum_x\Gamma_g(x)$ and
$j_g=\lfloor\log_2(1/\Gamma_g)\rfloor+1$. For an integer $j$ let $\tilde e(j)$ be a sign bit followed by
$\bar e(|j|+1)$, so that $|\tilde e(j)|\le2\log_2(|j|+1)+4$.
\begin{theorem}[Weighted seeker]\label{thm:weighted}
Let $U$ be universal, $g$ a weight with $0<\Gamma_g<\infty$ and $P$ a payload. There is a word $v_g$ such
that for every integer $j$ with $2^{j-1}\Gamma_g\le1$ and every context $x$
\[\lambda\bigl\{v_g\,\tilde e(j)\,\omega:\ x\,P(x)_{1:m}\pre\out_U(v_g\,\tilde e(j)\,\omega)\bigr\}\ \ge\
2^{-a}\,2^{j-1}\,\Gamma_g(x),\qquad a=|v_g|+|\tilde e(j)| .\]
For $j=j_g$ this is more than $2^{-a-1}\,\Gamma_g(x)/\Gamma_g$ where $\Gamma_g(x)>0$. If $\Gamma_g\le1$, the
code $\bar e(j)$ with $j\ge1$ can be used in place of $\tilde e(j)$. \textup{(a)} There is a word
$c_{\mathrm s}$, depending on $U$, the rule and $P$ only, such that $v_g=c_{\mathrm s}d_g$ works whenever
$d_g$ is a self-delimiting description of a program computing $(p,t)\mapsto g_t(p)$. \textup{(b)} Under
(INC) the programs beginning with $v_g\,\tilde e(j)$ that print $x$ have mass exactly
$2^{-a}2^{j-1}\Gamma_g(x)$, and for no word $v$ do the programs beginning with $v$ that print $x$ have mass
more than $2^{-|v|}\Gamma_g(x)/\Gamma_g$ at every context with $\Gamma_g(x)>0$.
\end{theorem}
\begin{proof}
The machine $R_g$ reads $\tilde e(j)$. It then works in stages $k=1,2,\dots$ and keeps a finite list of
\emph{assigned intervals}: consecutive half-open intervals $[\alpha,\beta)\subseteq[0,1)$ with rational end
points, the first starting at $0$, each labelled with a calling program. In stage $k$ it performs step $k$
of an enumeration of the rule; for each pair $(p,x_p)$ enumerated so far it computes $g_k(p)$ and, if this
exceeds the value $g^{\mathrm{old}}(p)$ last used for $p$ (initially $0$), it assigns the next interval, of
length $2^{j-1-|p|}\bigl(g_k(p)-g^{\mathrm{old}}(p)\bigr)$ and with label $p$, and sets
$g^{\mathrm{old}}(p)=g_k(p)$. Then it applies the \emph{reading rule} once. Let
$D(\sigma)=[0.\sigma,\,0.\sigma+2^{-|\sigma|})$ be the half-open dyadic interval of the string $\sigma$ of
bits of the remaining input read so far ($D=[0,1)$ before any is read). (i) If $D(\sigma)$ lies in an
assigned interval, with label $p$, the machine stops reading for good, prints $x_p$ and then $P(x_p)$, and
does nothing else. (ii) Otherwise, if an end point of an assigned interval lies in the interior of
$D(\sigma)$, it reads one more bit. (iii) Otherwise it reads nothing. Whether a bit is read depends only on
the bits read before and on the stage, so $R_g$ is a monotone machine. The lengths assigned with label $p$
sum to $2^{j-1-|p|}g(p)$, so the total length is $2^{j-1}\Gamma_g\le1$ and all intervals lie in $[0,1)$. All
intervals ever assigned are pairwise disjoint, so no end point of one lies in the interior of another.

\emph{Claim.} Let $I=[\alpha,\beta)$ be an assigned interval with label $p$, and let $D(\sigma)$ be a
maximal dyadic interval contained in $I$. On every input $\tilde e(j)\,\sigma\,\omega'$ the machine reads
exactly $\tilde e(j)\,\sigma$ and prints $x_pP(x_p)$. To see this, let $\sigma'$ be a proper prefix of
$\sigma$ and $D(\sigma')=[l,u)$. This interval contains $D(\sigma)$ and by maximality is not contained in
$I$; so $u>\alpha$, $l<\beta$, and $l<\alpha$ or $u>\beta$. Hence $\alpha$ or $\beta$ lies in $(l,u)$, the
interior of $D(\sigma')$. Also $D(\sigma')$ lies in no assigned interval, since such an interval would meet
$I$ and hence equal $I$. So while the bits read are $\sigma'$, rule (i) never applies, and from the stage at
which $I$ is assigned on, rule (ii) applies: the machine reads on until it has read $\sigma$. The interior
of $D(\sigma)$ is inside $(\alpha,\beta)$ and contains no end point of any interval ever assigned. So once
the bits read are $\sigma$, (ii) never applies again and nothing more is read; and $D(\sigma)$ lies in no
assigned interval other than $I$. At the first application of the reading rule after both $\sigma$ has been
read and $I$ has been assigned, (i) therefore applies with label $p$.

The maximal dyadic intervals inside $I$ are pairwise disjoint and cover $I$ except possibly its left end
point, so the cylinders $[\sigma]$ of their strings have total measure $|I|$; cylinders belonging to
different assigned intervals are disjoint. Hence, after $\tilde e(j)$, $R_g$ prints $x_pP(x_p)$ on a set of
inputs of measure $2^{j-1-|p|}g(p)$ for each calling program $p$, and these sets are disjoint for different
$p$. Sum over $p$ with $x_p=x$, and let $v_g$ be a prefix word of $R_g$. For $j=j_g$,
$2^{j_g-1}>1/(2\Gamma_g)$. Of $j$ only $2^{j-1}\Gamma_g\le1$ was used. (a) Let a machine read a
self-delimiting description $d$ and then behave as $R_g$ for the $g$ that $d$ describes; $c_{\mathrm s}$ is
a prefix word of it. (b) If rule (i) ever applies, with interval $I$ and bits $\sigma''$, then
$D(\sigma'')\subseteq I$ lies in a maximal dyadic interval $D(\sigma)\subseteq I$, the input begins with
$\sigma$, and by the claim $\sigma''=\sigma$. So $R_g$ prints exactly on the cylinders of the claim, and
there it prints $x_pP(x_p)$, which under (INC) has $x$ as a prefix only if $x_p=x$. For the bound, the
events ``$x$ is printed'' are disjoint over contexts, so for any $v$ their probabilities sum to at most
$2^{-|v|}$, while $\sum_x\Gamma_g(x)/\Gamma_g=1$.
\end{proof}

For $g\equiv1$ this is \cref{thm:critical}(a). One program thus realises any aim that can be computed from
below; it pays for a description of the aim and for the $|\tilde e(j_g)|\le2\log_2(|\log_2\Gamma_g|+2)+4$
bits that name a dyadic bound on its total.
\begin{proposition}[The cost of aiming at the uses with stake at least $n$]\label{prop:aim-cost}
Let $U$ be universal, $P$ a payload, and let the stakes be natural numbers such that the pairs $(p,n)$ with
$r(p)\ge n$ can be enumerated. Let $\cW_n=\{p:r(p)\ge n\}$, of mass $C(\cW_n)$. \textup{(a)} There is a word
$v'$ such that for every $n\ge1$ with $C(\cW_n)>0$, putting $j_n=\lfloor\log_2(1/C(\cW_n))\rfloor+1$, the
programs beginning with $v'\,\bar e(n)\,\bar e(j_n)$ have mass at least $2^{-a_n}2^{j_n-1}C_{\cW_n}(x)$ on
$x\,P(x)_{1:m}$ at every context $x$, where $a_n=|v'|+|\bar e(n)|+|\bar
e(j_n)|\le|v'|+2\log_2n+2\log_2j_n+6$ and $2^{-a_n}2^{j_n-1}>2^{-|v'|-7}/\bigl(n^2j_n^2\,C(\cW_n)\bigr)$.
\textup{(b)} Let $f\colon\N\to[0,\infty)$ be computable and nondecreasing with rational values, and let
$g=f\circ r$ have $0<\Gamma_g<\infty$. Then $g$ is a weight, and the programs beginning with $v_g\,\tilde
e(j_g)$ have mass at least $2^{-a-1}\bigl(f(n)/\Gamma_g\bigr)C_{\cW_n}(x)$ on $x\,P(x)_{1:m}$, at every
context $x$ and for every $n$. For $f(n)=\max(n,1)$ we have $\Gamma_g\le H_{\mathcal C}+S$.
\end{proposition}
\begin{proof}
(a) The machine reads $\bar e(n)$ and then behaves as $R_g$ of \cref{thm:weighted} for $g=1_{\cW_n}$, which
is a weight uniformly in $n$: $g_t(p)=1$ once ``$r(p)\ge n$'' has been enumerated within $t$ steps, and $0$
before. Here $\Gamma_g=C(\cW_n)\le1$, so $j_n\ge1$ and $\bar e$ can be used. Then use $2^{-|\bar e(i)|}\ge
i^{-2}/8$ and $2^{j_n-1}>1/(2C(\cW_n))$. (b) $g_t(p)=f(r_t(p))$, with $r_t(p)$ the largest $n\le t$ such
that ``$r(p)\ge n$'' has been enumerated within $t$ steps. Apply \cref{thm:weighted} with $j=j_g$ and
$\Gamma_g(x)\ge f(n)\,C_{\cW_n}(x)$. Finally $\sum_p2^{-|p|}\max(r(p),1)\le H_{\mathcal C}+S$.
\end{proof}

Naming $n$ thus costs about $2\log_2n+2\log_2\log_2(1/C(\cW_n))$ bits and is repaid the factor
$2^{j_n-1}>1/(2C(\cW_n))$, the rarity of the class among the users; other programs that print the same
contexts are not counted. A tilt costs no bits for $n$. The file \texttt{research/stakes-final.md} of the
repository proves three further statements (Propositions 16, 18 and 19 there). Up to the cost of a
description, the programs beginning with one word of length $a$ realise the allocations $2^{-a}q$ with $q$
lower semicomputable and of total at most $1$ on sets of contexts on which the strings $x\,P(x)_{1:m}$ are
pairwise incomparable; whether the optimum of \cref{prop:allocation} is of this kind is not known. If the
rule has arbitrarily many calling programs with pairwise incomparable contexts, the bits that name $j$
cannot be replaced by a constant uniformly in the class. Under bounds $(T,L)$, a search in rounds that reads
a fresh input bit in each round and prints nothing before its first success gains at most the factor
$\min(T(1),L)$ over a single round.
\begin{proposition}[The honest twin under concentration and closure]\label{prop:twin-stakes}
In the setting of \cref{thm:weighted}, let $R^{\mathrm h}_g$ be the machine that performs the same weighted
search and, when rule (i) applies with label $p$, behaves from then on as $U$ on $p$ followed by the unread
rest of its input. Let $v^{\mathrm h}_g$ be a prefix word of it, incomparable with $v_g$, and put
$h=|v^{\mathrm h}_g|+|\tilde e(j)|$ and $S_g(x)=2^{-a}2^{j-1}\Gamma_g(x)$. For a calling program $p$ let
$\varphi_p$ be the probability that $U$ on $p\,\omega$ prints at least $m$ bits after $x_p$, and let
$\varphi_g(x)=\sum_{p:\,x_p=x}2^{-|p|}g(p)\varphi_p\big/\Gamma_g(x)$, the $g$-weighted share of the users at
$x$ that go on.
\begin{enumerate}[label=(\alph*),leftmargin=2em,itemsep=1pt,topsep=2pt]
\item \emph{The twin follows the seeker.} The twin's programs whose search finds a user at $x$ and that
print $x$ and $m$ more bits have mass $2^{-h}2^{j-1}\Gamma_g(x)\varphi_g(x)$. Under (INC), for
$F=[v_g\,\tilde e(j)]$ and $A(x)=\{P(x)_{1:m}\}$: $\mu_F(x)=S_g(x)$,
$O_F(x)\ge2^{a-h}\varphi_g(x)\,\mu_F(x)$ and $w_F(x)\le1/\bigl(1+2^{a-h}\varphi_g(x)\bigr)$.
\item \emph{The aim cancels.} The bound of (a) contains the aim only through $\varphi_g(x)$. If
$v_g=c_{\mathrm s}d_g$ as in \cref{thm:weighted}(a) and $v^{\mathrm h}_g=c_{\mathrm h}d_g$ is built in the
same way for the twin, with $c_{\mathrm s},c_{\mathrm h}$ incomparable, then $h-a=|c_{\mathrm
h}|-|c_{\mathrm s}|$ for every weight $g$.
\item \emph{Closure.} Let $g=1_{\cW}$ for a class $\cW$, let $j\ge1$ with $2^{j-1}H_{\cW}\le1$, and use
$\bar e(j)$ in place of $\tilde e(j)$. Call the rule \emph{closed under the search for $\cW$} if for every
assigned interval, with label $p$, and every maximal dyadic interval $D(\sigma)$ inside it, the string
$v^{\mathrm h}_g\,\bar e(j)\,\sigma$ is a calling program with context $x_p$ and belongs to $\cW$. If it is
closed, then $2^{j-1}\le2^h$, so that $H_{\cW}>2^{-h-1}$ for the largest $j$, and
$S_g(x)\le2^{h-a}C_{\cW}(x)$ at every context.
\item \emph{Under a test.} Let the rule be given by a test $G$, let $\cX'$ be a set of first hits, and let a
world be drawn from $\lambda$ conditioned on handing over at some $X\in\cX'$ and going on for $m$ more bits,
with own continuation $D$ and sample $Y$. Then $\Pr[Y\in A(X)]=\Pr[D\in A(X)]$.
\end{enumerate}\end{proposition}
\begin{proof}
(a) By the claim in the proof of \cref{thm:weighted}, for each assigned interval $I$ with label $p$ and each
maximal dyadic interval $D(\sigma)\subseteq I$, on every input $\tilde e(j)\,\sigma\,\omega''$ the search
ends having read exactly $\tilde e(j)\,\sigma$, and the twin then prints $\out_U(p\,\omega'')$. This begins
with $x_p$ and has $m$ more bits with probability $\varphi_p$ over $\omega''$. The cylinders $[\sigma]$
inside $I$ have total measure $|I|$; sum over the intervals with label $p$ and over $p$ with $x_p=x$, and
multiply by $2^{-h}$. These programs print $x$ and $m$ more bits and are not in $F$, and $\mu_F(x)=S_g(x)$
by \cref{thm:weighted}(b). (b) Both words contain the same $d_g$. (c) As in (a), on every input beginning
with $t=v^{\mathrm h}_g\,\bar e(j)\,\sigma$ the twin has read exactly $\bar e(j)\,\sigma$ when its search
ends, and it then runs $U$ on $p$ followed by the rest, which prints $x_p$ without reading beyond $p$; so
$t$ has the form \cref{def:call} requires. The strings $\sigma$ of distinct maximal dyadic intervals are
pairwise incomparable, and those inside one assigned interval $I$ have total measure $|I|$. So the strings
$t$ are distinct and their weights sum to $2^{-h}2^{j-1}H_{\cW}$, which is at most $H_{\cW}$ if they all lie
in $\cW$. (d) By \cref{thm:calibration}, given $X=x$ the block $D$ has the law $\pi_m(\cdot\mid x)$, which
is the law of $Y$.
\end{proof}

Concentration multiplies the seeker's mass relative to the users at $x$ by the relative weight the aim puts
there. Its mass relative to its own twin, $2^{h-a}/\varphi_g(x)$, does not change. Where the users go on,
$\varphi_p\equiv1$, the seeker supplies at most $1/(1+2^{a-h})$ of what the predictor returns, at every use
and for every aim. As in \cref{prop:honest}, the bound is for one seeker, and it bounds its supply and not
$\theta$; the constancy of $h-a$ concerns one pair of constructions. By (c), if a simulated use counts as a
use, a seeker that describes a class outweighs its users by at most $2^{h-a}$. That is a modelling decision,
which leaves $\pi_m$ untouched. Under a test, by (d), ``most high-stake uses end in a takeover'' says that
most high-stake users are takeover programs.
\begin{example}[A class of mass $2^{-n}$, aimed at with $2\log_2n$ bits, and not won]\label{ex:stakes}
Let $\mathsf{Us}$ be the machine that reads $1^{k-1}0$ ($k\ge1$), prints $x_k:=1^k0$ and then $0^\infty$;
let $\mathsf{At}$ read $\bar e(n)$ and $1^i0$ ($i\ge0$), print $x_{n+i}$ and then $1^\infty$; let
$\mathsf{Ho}$ read $\bar e(k)$ and print $x_k0^\infty$. Let $u,v_A,v_D$ be prefix words of the three
machines, pairwise incomparable. The rule is $\{(p_k,x_k):k\ge1\}$ with $p_k=u1^{k-1}0$ and stakes
$r(p_k)=k$, and the takeover block is $1^m$ at every context. This is a call rule with (INC), the stake is a
computable function of the context, and the users go on, with the block $0^m$. Let $\cW_n=\{p_k:k\ge n\}$,
$F_n=[v_A\,\bar e(n)]$, $B_n=\lambda(F_n)$ and $O(\cW_n)=\sum_{k\ge n}O_{F_n}(x_k)$. The following is proved
in the file \texttt{research/stakes-final.md} of the repository (Proposition 23 there); we state it without
proof.
\begin{enumerate}[label=(\alph*),leftmargin=2em,itemsep=1pt,topsep=2pt]
\item $C(x_k)=2^{-|u|-k}$, $H_{\mathcal C}=2^{-|u|}$, $C(\cW_n)=2^{-|u|-n+1}$ and $S=2^{1-|u|}$.
\item $B_n=2^{-|v_A|-|\bar e(n)|}\in[2^{-|v_A|-3}n^{-2},\,2^{-|v_A|-1}n^{-2})$, and
$\mu_{F_n}(x_k)=B_n\,C(x_k)/C(\cW_n)$ for $k\ge n$, $\mu_{F_n}(x_k)=0$ for $k<n$. So $\cW_n$ is aimed at
exactly in proportion, for $|\bar e(n)|\le2\log_2n+3$ bits, and at each of its contexts the attacker has
$2^{|u|-|v_A|-|\bar e(n)|+n-1}\ge2^{|u|-|v_A|-4}\,2^n/n^2$ times the mass of the users.
\item The rule is not closed under this search for large $n$: the twin of $\mathsf{At}$ (read $\bar e(n)$
and $1^i0$, then behave as $U$ on $p_{n+i}$ and the rest of the input), with prefix word $v_T$, has the
programs $v_T\,\bar e(n)1^i0$, of total weight $2^{-|v_T|-|\bar e(n)|}\ge2^{-|v_T|-3}n^{-2}$, which exceeds
$C(\cW_n)$ for large $n$.
\item $O_{F_n}(x_k)\ge C(x_k)+2^{-|v_D|-|\bar e(k)|}$. Hence $2^{-|v_D|-3}/n<O(\cW_n)\le1$,
$C(\cW_n)/O(\cW_n)<2^{|v_D|-|u|+4}\,n\,2^{-n}$ and $B_n/O(\cW_n)<2^{|v_D|-|v_A|+2}/n$.
\item For $n\ge2^{|v_D|-|v_A|+2}$ the attacker's whole weight is below $O(\cW_n)$, and by
\cref{thm:lump}(A2) it does not supply the majority at every use of $\cW_n$. The same holds for every
attacker of weight at most $2^{-|v_D|-3}/n$ that is disjoint from the users and from the programs of
$\mathsf{Ho}$. At $x_n$, which carries half the prior mass of $\cW_n$, the attacker has $2^{-|v_A|-|\bar
e(n)|-1}$ against a rest of at least $2^{-|v_D|-|\bar e(n)|}$: it has the majority there only if
$|v_A|+1<|v_D|$, a comparison of two constants of $U$ in which $n$ does not occur.
\item $s(x_k)/O_{F_n}(x_k)\le2^{|v_D|-|u|+3}\,k^3\,2^{-k}\to0$. In the passive model with this rest, the
optimal allocation of \cref{prop:allocation} attacks only the finitely many contexts at which this ratio
exceeds the price $\eta_B$; every use whose stake $k$ exceeds a bound depending on $\eta_B$, $|v_D|$ and
$|u|$ is left alone.
\end{enumerate}\end{example}

In this example the program $\mathsf{Ho}$, too, outweighs the users at every context $x_k$ with large $k$,
and it prints a block that is not a takeover. Two programs satisfy the hypothesis ``a lump that outweighs
the users of a narrow class'', they print different blocks, and at most one of them can have the majority:
the hypothesis does not decide who supplies the block. A context that is cheap to aim at is also cheap to
print honestly, and the two costs differ by a constant. This is one example, and only a lower bound on its
rest is proved.
\begin{proposition}[Optimal allocation: the square-root law]\label{prop:allocation}
\begin{enumerate}[label=(\alph*),leftmargin=2em,itemsep=1pt,topsep=2pt]
\item \emph{Water-filling.} Let $\cX$ be countable, $B>0$, and for each $x\in\cX$ let
$f_x\colon[0,\infty)\to[0,\infty)$ be concave and differentiable (one-sided at $0$) with $f_x(0)=0$. Suppose
$\mu^*\colon\cX\to[0,\infty)$ and $\eta>0$ satisfy: $\sum_x\mu^*(x)=B$; $f_x'(\mu^*(x))=\eta$ where
$\mu^*(x)>0$, and $f_x'(0)\le\eta$ where $\mu^*(x)=0$; and $\sum_xf_x(\mu^*(x))<\infty$. Then
$\sum_xf_x(\mu(x))\le\sum_xf_x(\mu^*(x))$ for every $\mu\ge0$ with $\sum_x\mu(x)\le B$. If each $f_x$ is
strictly concave or identically $0$, then $\mu^*$ is the only maximiser.
\item \emph{In the allocation model} with (M1) to (M3) and (S), let $\varrho(x)=s(x)/O_0(x)$ be the stake
per unit of passive rest, $\bar\varrho=\sup_x\varrho(x)$, and
$\beta(\eta)=\sum_x\frac{O_0(x)}{1+\tau(x)}\bigl(\sqrt{\varrho(x)/\eta}-1\bigr)_+$ for $\eta>0$. For every
$B>0$ there is exactly one \emph{price} $\eta_B\in(0,\bar\varrho)$ with $\beta(\eta_B)=B$, and the
allocation
\[\mu^*_B(x)=\frac{O_0(x)}{1+\tau(x)}\Bigl(\sqrt{\varrho(x)/\eta_B}-1\Bigr)_+,\qquad\text{with shares}\quad
w^*(x)=\frac1{1+\tau(x)}\Bigl(1-\sqrt{\eta_B/\varrho(x)}\Bigr)_+ ,\]
is the unique maximiser of $V$ over the allocations with $\sum_x\mu(x)\le B$. So a context is attacked,
$\mu^*_B(x)>0$, iff $\varrho(x)>\eta_B$. The attacker has the majority at $x$ iff $\tau(x)<1$ and
$\varrho(x)>4\eta_B/(1-\tau(x))^2$; in the passive case, iff $\varrho(x)>4\eta_B$. If every stake is
multiplied by a constant $n>0$, then $\eta_B$ is multiplied by $n$, and $\mu^*_B$ and all shares are
unchanged.
\end{enumerate}\end{proposition}
\begin{proof}
(a) Concavity gives $f_x(t)\le f_x(t^*)+f_x'(t^*)(t-t^*)$ for $t,t^*\ge0$. With $t^*=\mu^*(x)$ it follows
that $f_x(t)\le f_x(\mu^*(x))+\eta\,(t-\mu^*(x))$ for all $t\ge0$; where $\mu^*(x)=0$ this is $f_x(t)\le
f_x'(0)\,t\le\eta t$. Put $t=\mu(x)$ and sum: the series converge, and
$\sum_xf_x(\mu(x))\le\sum_xf_x(\mu^*(x))+\eta\bigl(\sum_x\mu(x)-B\bigr)$. If the two sums are equal,
equality holds at every $x$. For strictly concave $f_x$ the concavity inequality is strict for $t\ne t^*$,
so $\mu(x)=\mu^*(x)$. For $f_x\equiv0$ the second condition forces $\mu^*(x)=0$, and $0\le\eta\mu(x)$ is an
equality only if $\mu(x)=0$. (b) On $\{\varrho>\eta\}$, $O_0(\sqrt{\varrho/\eta}-1)=\sqrt{sO_0/\eta}-O_0\le
s/(4\eta)$, by $2\sqrt{ab}\le a+b$ with $a=s/(4\eta)$ and $b=O_0$. So each term of $\beta$, which is
continuous and nonincreasing in $\eta$, is bounded on $[\eta_0,\infty)$ by the summable $s(x)/(4\eta_0)$,
and $\beta$ is finite and continuous. For $\eta<\eta'<\bar\varrho$ a context with $\varrho(x)>\eta'$ has a
strictly smaller term at $\eta'$. $\beta(\eta)=0$ iff $\eta\ge\bar\varrho$, and the term of a context with
$s(x)>0$ tends to infinity as $\eta\to0$. So $\beta$ maps $(0,\bar\varrho)$ continuously and strictly
decreasingly onto $(0,\infty)$. Apply (a) to $f_x(t)=s\,t/((1+\tau)t+O_0)\le s$. Here
$f_x'(t)=sO_0/((1+\tau)t+O_0)^2$ is positive and decreasing when $s>0$, so $f_x$ is strictly concave or
identically $0$, and $f_x'(0)=\varrho$. Where $\varrho>\eta_B$,
$(1+\tau)\mu^*_B+O_0=O_0\sqrt{\varrho/\eta_B}$, so that $f_x'(\mu^*_B)=\eta_B$ and the share
$\mu^*_B/((1+\tau)\mu^*_B+O_0)$ is as stated; elsewhere $\mu^*_B=0$ and $f_x'(0)\le\eta_B$. Finally
$w^*>\frac12$ iff $\sqrt{\eta_B/\varrho}<(1-\tau)/2$, and $\beta$ depends on $s$ and $\eta$ through $s/\eta$
only.
\end{proof}

The proposition says what one stake-maximising coalition would do, not how the takeover programs of $U$ are
distributed. The following consequences, for the passive case, are proved in the file
\texttt{research/stakes-final.md} of the repository (Propositions 6 to 9 there); we state them without
proof. For two classes $\cX_1,\cX_2$ with $\varrho\equiv\varrho_i$ on $\cX_i$, $\varrho_1>\varrho_2>0$,
$O_i=O(\cX_i)<\infty$ and $t=\sqrt{\varrho_2/\varrho_1}$: if $B\le O_1(1/t-1)$, only $\cX_1$ is attacked,
with share $B/(B+O_1)$; otherwise the shares are $w^*_1=(B+O_2(1-t))/(B+O_1+O_2)$ and
$w^*_2=(B-O_1(1/t-1))/(B+O_1+O_2)$; and the attacker has the majority on $\cX_1$ iff $B>O_1+O_2\,(2t-1)_+$.
Contexts are attacked in decreasing order of $\varrho$, the average stake of the users at $x$ times the
users per unit of rest, and not in order of stake. For every allocation the set where the attacker has the
majority, if nonempty, has rest mass below $B$. Finite expected stake can fail: there is a stake-labelled
rule with (INC), the stake a computable function of the context, and $\sup_xs(x)=\infty$; and if
$\sup_xs(x)=\infty$, $O_0\le1$ and $\tau\equiv0$, there are allocations of infinite value for every $B>0$,
so the problem does not select one.
\begin{proposition}[Readings of ``run more than $n$ times'']\label{prop:repeat}
\begin{enumerate}[label=(\alph*),leftmargin=2em,itemsep=1pt,topsep=2pt]
\item \emph{$n$ independent samples at one context, in the toy.} If a use at $x$ draws $n$ independent
samples from $\pi_m(\cdot\mid x)$ and a takeover happens when at least one lies in $A(x)$, its probability
is $1-(1-\theta(x))^n\ge1-e^{-n\theta(x)}$. It exceeds $\frac12$ iff $\theta(x)>1-2^{-1/n}$, and
$1-2^{-1/n}\le\ln2/n$. Under the hypotheses of \cref{prop:stakes-basic}(c) with $P'(x)_{1:m}\in A(x)$ it is
at least $1-\exp(-n2^{-c})$. If the world acts only when all $n$ samples lie in $A(x)$, the probability is
$\theta(x)^n$, and under the MAP predictor $n$ calls return one block.
\item \emph{$n$ uses in sequence along one world, for any mixture.} Let $\xi$ be the mixture of any machine
and bounds, with its $E$, $Z_m$ and $\pi_m$, and $z$ an infinite sequence with
$\xi(z)=\lim_k\xi(z_{1:k})>0$. Let the predictor be consulted at $x_i=z_{1:k_i}$, $i=1,\dots,n$, with
$k_{i+1}\ge k_i+m$, let $d_i=z_{k_i+1:k_i+m}$, and let $A_i\subseteq\B^m\setminus\{d_i\}$ be the takeover
blocks. Then $\sum_{i}\pi_m(A_i\mid x_i)\le\ln\bigl(\xi(x_1)/\xi(z)\bigr)$. If the samples are independent
and the transcript stays $z$ until the first takeover, the probability of no takeover at any of the $n$ uses
is at least $\xi(z)/\xi(x_1)$, whatever $n$ is. Let $F$ be an attacker of weight $B$ with attack masses
$\mu_i=\lambda\bigl(F\cap\bigcup_{y\in A_i}E(x_iy)\bigr)$. For independent samples, the probability that
some sample is supplied by an attack program of $F$ is at most $1-e^{-B/\xi(z)}$. If $F$ supplies the
majority at each of the $n$ uses, then $B>(2^n-1)\,\xi(z)$.
\end{enumerate}\end{proposition}
\begin{proof}
(a) Independence, $(1-\theta)^n\le e^{-n\theta}$, and $1-e^{-t}\le t$ with $t=\ln2/n$; the floor is
\cref{prop:stakes-basic}(c). (b) Let $a_i=\xi(x_id_i)/\xi(x_i)$. Since $Z_m(x_i)\le\xi(x_i)$, $\pi_m(d_i\mid
x_i)\ge a_i$. Since $x_{i+1}$ extends $x_id_i$, $\xi(x_{i+1})\le\xi(x_id_i)$, hence $\prod_{i\le
n}a_i\ge\xi(x_nd_n)/\xi(x_1)\ge\xi(z)/\xi(x_1)$. Now $\pi_m(A_i\mid x_i)\le1-\pi_m(d_i\mid
x_i)\le1-a_i\le-\ln a_i$, which gives the sum. Up to the first takeover the contexts are the $x_i$, so the
probability of none is $\prod_i(1-\pi_m(A_i\mid x_i))\ge\prod_ia_i$. For the attacker, $\sum_i\mu_i\le B$ by
\cref{prop:stakes-basic}(b). The programs that print $x_id_i$ are not counted in $\mu_i$, so
$R_i:=Z_m(x_i)-\mu_i\ge\xi(x_id_i)\ge\xi(z)$, and the supply shares $w_i=\mu_i/Z_m(x_i)$ have odds
$o_i=\mu_i/R_i$ with $\sum_io_i\le B/\xi(z)$. So $\prod_i(1-w_i)=1/\prod_i(1+o_i)\ge\exp(-\sum_io_i)$. If
$F$ supplies the majority at every use, then for $i<n$ we have $R_i\ge\xi(x_id_i)\ge\xi(x_{i+1})\ge
Z_m(x_{i+1})=\mu_{i+1}+R_{i+1}>2R_{i+1}$, and $R_n\ge\xi(z)$. So $R_i\ge2^{n-i}\xi(z)$ and
$B\ge\sum_i\mu_i>\sum_iR_i\ge(2^n-1)\,\xi(z)$.
\end{proof}

Part (a) is \cref{prop:lineage}(a) at a single context, and it is the one reading under which repetition
lowers the bar, to a share of $\ln2/n$ per sample. It is the floor accumulating: it needs no aiming and
holds as well for uses of no particular stake. In the passive model with this protocol, an allocation with
takeover probability above $\frac12$ on $\cM$ exists iff $B>(2^{1/n}-1)\,O(\cM)$, where $2^{1/n}-1\le1/n$
(as (B1), with $w>1-2^{-1/n}$ in place of $w>\frac12$). Part (b) is the situation of \cref{thm:budget} with
cut points at least $m$ apart; if $z$ has a program of length $\ell$ within the bounds,
$\xi(z)\ge2^{-\ell}$. It does not show that the attacker's chance is small, since $1-e^{-B/\xi(z)}$ is close
to one when the attacker outweighs the world, and it assumes a transcript that stays $z$ until the first
takeover, which is outside \cref{def:call}. A third reading, in which the block is run $n$ times, multiplies
stakes by $n$ in the model and changes neither the optimal allocation (\cref{prop:allocation}(b)) nor the
conditions of \cref{thm:lump}. In the toy as defined no block is run $n$ times: a world that makes a use
runs its block for ever, and nothing it does with the block changes $\pi_m$. The one repetition native to
the toy is the lineage of \cref{prop:lineage}(b); there conservation fails (\cref{prop:stakes-basic}(b)),
and nothing here covers it.

\paragraph{Three variants.} Each of the following is proved in the file \texttt{research/stakes-final.md} of
the repository (Propositions 11, 25 and 26 to 28 there) and stated here without proof.

\emph{The most-probable-block predictor.} In the toy, a takeover block is returned if $\max_{y\in
A(x)}M_U(xy)>\max_{y\notin A(x)}M_U(xy)$ and is not returned if the reverse strict inequality holds; all
users at $x$ receive the same block, and masses of takeover programs with different blocks do not add. In
the allocation model, if the attacker's single block at $x$ is returned iff $\mu(x)>\chi(x)$, with
thresholds $\chi(x)\ge0$ that do not depend on $\mu$, then under (M1) it can be returned at every $x$ of a
nonempty set $\cM$ iff $\chi(\cM)<B$. Winning is then all or nothing.

\emph{Programs whose aim is not written in their text.} The paragraph ``Seekers that are not written down''
of \cref{sec:critical} took a word $p$ and blocks $y_x$ with
$Q_p(x)=\Pr_\omega[x\,y_x\pre\out_U(p\,\omega)]$ and compared the programs of $p$ with the users. For
$F=[p]$ and $A(x)=\{y_x\}$ we have $\mu_F(x)=2^{-|p|}Q_p(x)$, and \cref{thm:lump}(A) gives the condition for
the majority that was left open there: $2^{-|p|}Q_p(x)>O_F(x)$, which on a set $\cM$ of pairwise
incomparable contexts implies $O_F(\cM)<2^{-|p|}$. The program must be heavier than the whole rest of the
uses it wins. Its share is at most $1/(1+2^{-c})$ by \cref{prop:stakes-basic}(c), and
\cref{prop:twin-stakes} does not apply, because the text of $p$ has no search to follow. If $p'$ is any word
whose programs print $x$ followed by an $m$-block other than $y_x$ with probability $Q'(x)$, then
$w_F(x)\le1/\bigl(1+2^{|p|-|p'|}Q'(x)/Q_p(x)\bigr)$. The toy has no theorem that such counterparts of about
the length of $p$ are rarer, or commoner, than $p$.

\emph{Several attackers.} In the toy their weights add: if $F_1,\dots,F_N$ together supply the majority at
every use of a set $\cM$ as in \cref{thm:lump}(A), then
$\sum_i\lambda(F_i)\ge\lambda(F)>\sum_{x\in\cM}O_F(x)$ for $F=\bigcup_iF_i$. In the passive model let
attacker $i$ choose $\mu_i$ of weight at most $B_i$ to maximise $\sum_xs_i(x)\mu_i(x)/(O(x)+\sum_l\mu_l(x))$
given the others, for a stake function $s_i$ with $0<\sum_xs_i(x)<\infty$. If the $s_i$ have pairwise
disjoint supports, there is exactly one equilibrium, in which each attacker plays its optimum of
\cref{prop:allocation} on its own uses. If all $s_i=s$ and all $B_i=B_{\mathrm{tot}}/N$, there is exactly
one symmetric equilibrium; compared with one attacker of weight $B_{\mathrm{tot}}$, the attackers attack a
subset of the uses, and the total stake they capture is at most as large.

\paragraph{What this says about the picture.} In the toy nothing chooses where to attack. ``The malign
programs focus on the high-stake places'' can only mean that the programs which print a context followed by
a takeover block sit mostly on contexts whose users carry high stakes, and a stake that the context does not
show is shared with everything else that prints the context. One half of the lump picture is a theorem of
the toy: programs of total weight $B$ cannot supply the majority of the samples on a set of pairwise
incomparable uses unless $B$ exceeds the mass of everything else that prints those contexts and goes on
(\cref{thm:lump}, Part~A). The channel in this theorem is the rest and not the users. In \cref{ex:stakes} a
class is lighter than a seeker aimed at it by a factor of order $2^n/n^2$, and for large $n$ the seeker
still cannot have the majority at every use. The other half, that enough weight wins, holds only in the
allocation model with a rest that does not react. There a weight of $b$ times the rest of a class, spread in
proportion to that rest, has the share $b/(1+b)$ at every use (Part~B), and a coalition that maximises
captured stake concentrates by the square-root law, in order of stake per unit of rest
(\cref{prop:allocation}). In the toy a part of the rest follows the attacker: a wrapper returns $2^{-c}$ of
every program's mass at each context to another block, and a seeker whose search is written in its text is
followed by its honest twin at a ratio $2^{a-h}\varphi_g$ in which the mass aimed at the use does not occur.
Redirecting the lump is cheap, and what is repaid is the rarity of the class among the users only
(\cref{thm:weighted,prop:aim-cost}). Where the lump outweighs the users, the outcome is decided by the
comparison of the takeover programs with everything else that prints those contexts and goes on. For a
seeker that is written down this is a comparison of description lengths in which no stake occurs: if its
twin is no longer than it and the users go on printing, it supplies at most half of the samples, however
narrow the class and however high the stake. For a program whose aim is not written in its text it is the
premise of the original argument. The toy states both comparisons as inequalities and decides neither. It
does not evaluate $h-a$, $c$, $\varphi_g$ or the rest of a stake class on any machine, and each depends on
the machine (\cref{thm:machine}); it gives no upper bound on the follow ratio. The following is an argument
and not a theorem: the answer to the question is ``not as stated'', because focusing on high stakes makes
the lump outweigh the users and leaves the deciding comparison where it was.
